\section{Introduction}
\label{sec:introduction}

To audit any computational process, we need robust and well-founded notions of \emph{provenance} to track how
data are used. This allows us to answer questions like ``Where did these data come from?'', ``Why are these
data in the output?'' and ``How were these data computed?''. Provenance tracking has a wide range of
applications, from debugging and program comprehension~\cite{buneman95,cheney07} to improving reproducibility
and transparency in scientific workflows~\cite{kontogiannis08}. \emph{Program slicing}, first proposed
by~\citet{weiser81}, is a collection of techniques for provenance tracking that attempts to take a run of a
program and areas of interest in the output, and turn them into the subset of the input and the program that
were responsible for generating those specific outputs.

Underlying all of these questions is a more basic one: how does the output of a program respond when
its input changes? For programs that compute with real numbers, calculus provides a canonical answer: the
derivative, a linear map between tangent spaces describing how the output varies as the input varies, with
\emph{automatic differentiation} (AD) being a computational technique for computing such derivatives alongside
the program itself~\cite{siskind08,elliott18,vakar22}. For data provenance and dependency analysis, where the
inputs are database tuples or elements of a data structure, the question is often framed in a more qualitative
way: if we perturb the input at a given position, can the output change at all? Posing this question one input
at a time yields a vector of Boolean partial ``derivatives'', and for a program with several outputs, a
\emph{Boolean Jacobian}: a matrix whose entry at $(j, i)$ records whether output position $j$ may depend on
input position $i$. Consider multiplication: at the point $(x, y)$, the usual partial derivative $\partial(x
\cdot y)/\partial x$ is $y$, and its Boolean counterpart records that the output can vary with $x$ only when
$y \neq 0$. Like their numerical counterparts, Boolean Jacobians compose by matrix multiplication, with
conjunction and disjunction playing the role of multiplication and addition. This paper presents an approach
to dynamic dependency analysis where such Jacobians are evaluated alongside the program, forwards or
backwards, in the manner of forward- and reverse-mode AD, with the two directions related by matrix
transposition.

Existing approaches to program slicing are often tied to particular programming languages or implementations.
In this paper we develop this analogy into a general categorical approach to data provenance, in
which the derivative of a program is a linear map between spaces of approximations of its input and output,
with scalars drawn from a commutative semiring of dependency information. The choice of semiring determines
the analysis: the reals recover AD itself, while over distributive lattices the derivative and its transpose
form a \emph{conjugate pair}, propagating dependency information forwards and backwards; when the lattices are
Boolean algebras, the same maps can instead be presented as an adjunction, recovering the \GPS of Perera and
collaborators. Our categorical approach should provide a suitable setting for enabling ``automatic''
data provenance for a variety of programming languages, and is easily configured to use alternative
approximation strategies, including quantitative forms of slicing.

\subsection{Dependencies as Derivatives}
\label{sec:introduction:galois-slicing}

Consider a single run of a program on a particular input. Each part of the output is computed from
certain parts of the input, and recording these facts gives the \emph{dependency relation} of the run,
relating each output position to the input positions it depends on. This section works through a small example
showing how dependency relations behave as derivatives: they have a tabular representation as Boolean
matrices, compose using matrix multiplication, and can be evaluated forwards or backwards.

\begin{example}
  \label{ex:introduction-example}
  The following program is written in Haskell syntax \cite{haskell}, using a list comprehension to filter a list of pairs of labels and numbers to those numbers with a given label, and then computing the sum of the numbers:
  \begin{displaymath}
    \begin{array}{l}
      \mathrm{query} :: \mathrm{Label} \to [(\mathrm{Label}, \mathrm{Int})] \to \mathrm{Int} \\
      \mathrm{query}\,l\,\mathit{db} = \mathrm{sum}\,[ n \mid (l',n) \leftarrow \mathit{db}, l \equiv l' ]
    \end{array}
  \end{displaymath}
  With $\mathit{db} = [(\mathsf{a}, 0), (\mathsf{b}, 1), (\mathsf{a}, 1)]$, we will have $\mathit{query}\,\mathsf{a}\,\mathit{db}$ and $\mathit{query}\,\mathsf{b}\,\mathit{db}$ both evaluating to $1$.

  Suppose we are interested in how the output depends on the numerical parts of the input, for the
  query parameters $l = \mathsf{a}$ and $l = \mathsf{b}$. The three numbers in $\mathit{db}$
  occupy three \emph{positions}, and the output a single position. In the run with $l = \mathsf{a}$, the
  output is computed from the numbers at the first and third positions; in the run with $l = \mathsf{b}$,
  from the number at the second. We depict each input position pointing to the output positions that depend
  on it:
      \begin{center}
    \begin{tikzpicture}[every node/.style={font=\small}]
      % query a
      \node (la) at (0,1.4) {$\mathit{query}\;\mathsf{a}$};
      \node (a1) at (0,0.8) {$(\mathsf{a},0)$};
      \node (a2) at (0,0) {$(\mathsf{b},1)$};
      \node (a3) at (0,-0.8) {$(\mathsf{a},1)$};
      \node (ao) at (2.4,0) {$1$};
      \draw[->] (a1) -- (ao);
      \draw[->] (a3) -- (ao);
      % query b
      \node (lb) at (6,1.4) {$\mathit{query}\;\mathsf{b}$};
      \node (b1) at (6,0.8) {$(\mathsf{a},0)$};
      \node (b2) at (6,0) {$(\mathsf{b},1)$};
      \node (b3) at (6,-0.8) {$(\mathsf{a},1)$};
      \node (bo) at (8.4,0) {$1$};
      \draw[->] (b2) -- (bo);
    \end{tikzpicture}
  \end{center}
    Tabulating each relation, with a row for each output position and a column for each input position,
  gives the \emph{Jacobian} of the run:
    \begin{displaymath}
    \partial(\mathit{query}\,\mathsf{a}\,\mathit{db}) =
    \begin{pmatrix} 1 & 0 & 1 \end{pmatrix}
    \qquad
    \partial(\mathit{query}\,\mathsf{b}\,\mathit{db}) =
    \begin{pmatrix} 0 & 1 & 0 \end{pmatrix}
  \end{displaymath}
    As in differential calculus, the Jacobian is taken at a point: the two runs of the same program have
  different Jacobians. And like their numerical counterparts, these matrices compose: $\mathit{query}\,l$ is a
  sum applied to a selection, and the Jacobian of the run is the product of the Jacobians of those two stages,
  a chain rule for dependencies.

  In setting up the positions, we chose to make only the numbers perturbable, keeping the labels and
  the structure of the list itself fixed. Other choices are also useful, and one of the aims of this work is
  to clarify how to choose a dependency structure appropriate for different tasks by selecting an appropriate
  semiring of dependency information. We elaborate on this further in
  \secref{semantic-gps}.

  The Jacobian acting directly evaluates dependencies \emph{forwards}: given a selection of input
  positions, it returns the output positions that depend on them. A selection of positions can be displayed as
  an \emph{approximation} of the underlying value, with the numbers at unselected positions replaced by
  $\bot$. Writing $\partial(\mathit{query}\,l\,\mathit{db})_f$ for this forward action, we have
  $\partial(\mathit{query}\,\mathsf{a}\,\mathit{db})_f([(\mathsf{a},0),(\mathsf{b},\bot),(\mathsf{a},\bot)]) =
  1$, since the output depends on the entry $(\mathsf{a},0)$, whereas
  $\partial(\mathit{query}\,\mathsf{b}\,\mathit{db})_f$ applied to the same selection returns $\bot$: that run
  consulted no $\mathsf{a}$-labelled entries.

  The transpose of the Jacobian evaluates dependencies \emph{backwards}: given a selection of
  output positions, it returns the input positions on which they depend. Writing
  $\partial(\mathit{query}\,l\,\mathit{db})_r$ for this reverse action, the fully approximated output $\bot$
  requires none of the input:
  \begin{displaymath}
    \partial (\mathit{query}\,l\,\mathit{db})_r(\bot) = [(\mathsf{a},\bot), (\mathsf{b}, \bot), (\mathsf{a}, \bot)]
  \end{displaymath}
  for both $l = \mathsf{a}$ and $l = \mathsf{b}$. If instead we retain the output ``$1$''
  unapproximated, then the two
  query runs' backwards actions return different results:
  \begin{displaymath}
    \begin{array}{l}
      \partial (\mathit{query}\,\mathsf{a}\,\mathit{db})_r(1) = [(\mathsf{a},0), (\mathsf{b},\bot), (\mathsf{a},1)] \\
      \partial (\mathit{query}\,\mathsf{b}\,\mathit{db})_r(1) = [(\mathsf{a},\bot), (\mathsf{b},1), (\mathsf{a},\bot)]
    \end{array}
  \end{displaymath}
  Pieces of the input that were {\em not} used are replaced by $\bot$. As we expect, the run of the query with label $\mathsf{a}$ depends on the entries in the database labelled with $\mathsf{a}$, and likewise for the run with label $\mathsf{b}$.

  In a simple query like this, it is easy to work out the dependency relationship between the input and output. However, the benefit of language-based approaches is that they are {\em automatic} for all programs, no matter how complex the relationship between input and output. Moreover, by changing what we mean by ``approximation'' we can compute a range of different information about a program.
\end{example}

\subsection{Galois Slicing and Automatic Differentiation}

The backward maps of the previous section compute a form of \emph{data provenance}: for each part of
the output, the part of the input needed to explain it. Dependency information of this kind is the common
currency of a range of existing techniques, from program slicing~\cite{weiser81} to provenance for database
queries~\cite{buneman95,cheney07}, including semiring provenance~\cite{green07}. Closest
to our setting is the \GPS of Perera and
collaborators~\cite{perera12a,perera16d,ricciotti17}, which organises the forward and backward maps of a run
into a Galois connection between lattices of approximations.

Many forms of dynamic provenance analysis, \GPS included, work by recording a trace of each
execution and computing the backwards analysis by re-running over the trace. A denotational account would
instead bake the analysis into the semantics of the program itself, rather than provide it as a separately
defined ``backwards evaluation'' operation. 
The differential reading of dependency analysis, developed in \secref{approx-as-tangents}, points to
a way to obtain such an account, by analogy with \emph{automatic
differentiation}~\cite{siskind08,elliott18,vakar22}; let us make the correspondence explicit.

\begin{itemize}[leftmargin=\enummargin]
\item For dependency analysis, every value has an associated join-semilattice of \emph{approximations}, with least element $\bot$ (the empty selection). For differentiable programs, every point has an associated vector space of {\em tangents}.
\item For dependency analysis, every program has an associated forward approximation map that takes approximations of the input to the approximations of the output that depend on them. This map {\em preserves joins and $\bot$}. For differentiable programs, every program has a forward derivative that takes tangents of the input to tangents of the output. The forward derivative map is {\em linear}, so it preserves addition of tangents and the zero tangent.
\item For dependency analysis, every program has an associated backward approximation map that takes approximations of the output back to least approximations of the input that they depend on. This map also {\em preserves joins and $\bot$}. For differentiable programs, every program has a reverse derivative that takes \emph{cotangents} of the output (linear maps from tangents to scalars) to cotangents of the input. This map is again {\em linear}.
\item In both settings, the forward and backward maps are related by being each other's transpose.
The transpose makes sense because the spaces are \emph{self-dual}: tangent vectors pair with tangent vectors
via the dot product, identifying cotangents with tangents, and selections of positions pair by the same dot
product formula, with joins and meets in place of sums and products.
\end{itemize}

Given this close connection between dependency analysis and differentiable programming, we can take structures intended for modelling automatic differentiation, such as V{\'a}k{\'a}r's CHAD framework and use them to model dependency analysis. This will enable us to generalise and expand the scope of automatic differentiation to act as a foundation for data provenance in a wider range of computational settings.

\subsection{Outline and Contributions}

%Our primary contribution in this article is to elucidate \GPS by relating it to differentiable programming and automatic differentiation. In doing so, we aim to expand the applicability of program slicing, and to potentially transfer efficient automatic differentiation implementation techniques from automatic differentiation to program slicing and provenance tracking.

Our main contribution is to show that data provenance can be understood as differentiation over a
commutative semiring of dependency information. In \secref{approx-as-tangents} we develop the first-order
picture, for programs over tuples of atomic data: derivatives as matrices over a semiring $S$, with the
transpose relating forward and backward dependency analysis, and with successively stronger conditions on $S$ yielding
join-preserving maps, conjugate pairs and, for Boolean algebras, the Galois connections of \GPS; ordinary AD
occupies the case $S = \mathbb{R}$. We also relate this picture to \citet{berry79}'s \emph{stable domain
theory}, where stable functions provide a kind of semantic provenance analysis.

The following table summarises the correspondence developed over the course of the paper:
\vspace{1.5mm}
\begin{center}
  \small
  \begin{tabular}{lll}
    \toprule
    & \textbf{Differentiable programming} & \textbf{Dependency analysis} \\
    \midrule
    scalars & real numbers & commutative semiring $S$ \\
    space at a point & vector space of tangents & semimodule of approximations over $S$ \\
    identification with the dual & finite dimensionality & chosen self-duality \\
    derivative at a point & Jacobian matrix & $S$-weighted relation as matrix over $S$ \\
    chain rule & matrix multiplication & matrix multiplication over $S$ \\
    forward mode & push tangents forward & propagate dependencies forward \\
    reverse mode & pull cotangents back & propagate dependencies backwards \\
    forward/backward relation & transpose & transpose via self-dualities \\
    \bottomrule
  \end{tabular}
\end{center}
\vspace{1.5mm}

In \secref{models-of-total-gps} we extend to sum types using the category of families construction and,
following V{\'a}k{\'a}r \etal{}'s CHAD framework~\cite{vakar22,nunes2023}, build models of a higher-order
language in which every program is interpreted together with its family of derivatives. We apply this to a
concrete higher-order language in \secref{language} and demonstrate the model on a number of examples, showing
how parameterising on $S$ controls the dependency structure associated with data, something that was
``hard-coded'' in previous presentations of \GPS. We prove two correctness properties in
\secref{definability}: programs of first-order type denote matrices in the sense of
\secref{approx-as-tangents}, and the higher-order interpretation agrees with the standard set-theoretic
semantics. \secref{related-work} and \secref{conclusion} discuss additional related and future work.

% \begin{enumerate}[leftmargin=\enummargin]
% \item We explain how the CHAD framework of Vákár et al. can be adapted to give a general categorical framework for \GPS.
% \item Using our adapted CHAD framework, we can explain how type structure can be used to control the approximation lattices associated with data points, something that was ``hard coded'' in previous presentations of \GPS.
% \item With the benefit of our abstract setting, we can relate \GPS to other parts of denotational semantics. In particular, we show that there is a close connection with Stable Domain Theory, a proposed solution to capturing sequentiality by recording more intensional information about programs' sensitivity to approximations.
% \end{enumerate}

We have formalised our major results in Agda, resulting in an executable implementation built directly from the categorical constructions that we have used to compute the examples in \secref{language:examples}. Please consult the file \texttt{everything.agda} in the supplementary material.

% \begin{enumerate}
% \item We explain how the CHAD framework of Vákár et al. can be adapted to give a general categorical framework for \GPS.
% \item Using our adapted CHAD framework, we can explain how type structure can be used to control the approximation lattices associated with data points, something that was ``hard coded'' in previous presentations of \GPS.
% \item With the benefit of our abstract setting, we can relate \GPS to other parts of denotational semantics. In particular, we show that there is a close connection with Stable Domain Theory, a proposed solution to capturing sequentiality by recording more intensional information about programs' sensitivity to approximations.

% We have \anonymise{\href{https://github.com/bobatkey/approx-diff}{formalised this work in Agda}}{formalised this work in Agda}. Not only does this mean that the constructions and proofs have been checked, but also the construction is executable and can be run on small examples. We are also planning to use Agda's JavaScript backend to produce a demo.
